\Chapter{17}

\Verse{1}
{
    \zA{话说秦钟既死，宝玉痛哭不止，李贵等好容易劝解半日方住，归时还带馀哀。 贾母帮了几十两银子，外又另备奠仪，宝玉去吊祭。 七日后便送殡掩埋了，别无记 述。
        只有宝玉日日感悼，思念不已，然亦无可如何了。 又不知过了几时才罢。 这日贾珍等来回贾政：“园内工程俱已告竣，大老爷已瞧过了，只等老爷瞧了，
        或有不妥之处，再行改造，好题匾额对联。” 贾政听了，沉思一会，说道：“这匾 对倒是一件难事。 论礼该请贵妃赐题才是，然贵妃若不亲观其景，亦难悬拟。}
}
{
    \eA{Ch'in Chung, to resume our story, departed this life, and Pao-yü went on so
        unceasingly in his bitter lamentations, that Li Kuei and the other servants
        had, for ever so long, an arduous task in trying to comfort him before he
        desisted; but on his return home he was still exceedingly disconsolate.
        Dowager lady Chia afforded monetary assistance to the amount of several tens
        of taels;}
}
{
}

\Verse{2}
{
    \zA{若直 待贵妃游幸时再行请题，若大景致，若干亭榭，无字标题，任是花柳山水，也断不 能生色。” 众清客在旁笑答道：“老世翁所见极是。
        如今我们有个主意：各处匾对 断不可少，亦断不可定。 如今且按其景致，或两字、三字、四字，虚合其意拟了来，
        暂且做出灯匾对联悬了，待贵妃游幸时，再请定名，岂不两全？” 贾政听了道：“所 见不差。 我们今日且看看去，只管题了，若妥便用；
        若不妥，将雨村请来，令他再 拟。”}
}
{
    \eA{and exclusive of this, she had sacrificial presents likewise got ready.
        Pao-yü went and paid a visit of condolence to the family, and after seven
        days the funeral and burial took place, but there are no particulars about
        them which could be put on record. Pao-yü, however, continued to mourn (his
        friend) from day to day, and was incessant in his remembrance of him, but
        there was likewise no help for it.}
}
{
}

\Verse{3}
{
    \zA{众人笑道：“老爷今日一拟定佳，何必又待雨村。” 贾政笑道：“你们不知： 我自幼于花鸟山水题咏上就平平的，如今上了年纪，且案牍劳烦，于这怡情悦性的
        文章更生疏了。 便拟出来，也不免迂腐，反使花柳园亭因而减色，转没意思。” 众 清客道：“这也无妨。 我们大家看了公拟，各举所长，优则存之，劣则删之，未为
        不可。” 贾政道：“此论极是。 且喜今日天气和暖，大家去逛逛。” 说着，起身引 众人前往。 贾珍先去园中知会。}
}
{
    \eA{Neither is it known after how many days he got over his grief. On this day,
        Chia Chen and the others came to tell Chia Cheng that the works in the
        garden had all been reported as completed, and that Mr. Chia She had already
        inspected them. "It only remains," (they said), "for you, sir, to see them;}
}
{
}

\Verse{4}
{
    \zA{可巧近日宝玉因思念秦钟，忧伤不已，贾母常命人带他到新园子里来玩耍。 此 时也才进去，忽见贾珍来了，和他笑道：“你还不快出去呢，一会子老爷就来了。”
        宝玉听了，带着奶娘小厮们，一溜烟跑出园来。 方转过弯，顶头看见贾政引着众客 来了，躲之不及，只得一旁站住。 贾政近来闻得代儒称赞他专能对对，虽不喜读书，
        却有些歪才，所以此时便命他跟入园中，意欲试他一试。 宝玉未知何意，只得随往。 刚至园门，只见贾珍带领许多执事人旁边侍立。}
}
{
    \eA{and should there possibly be anything which is not proper, steps will be at
        once taken to effect the alterations, so that the tablets and scrolls may
        conveniently be written." After Chia Cheng had listened to these words, he
        pondered for a while. "These tablets and scrolls," he remarked, "present
        however a difficult task.}
}
{
}

\Verse{5}
{
    \zA{贾政道：“你且把园门关上， 我们先瞧外面，再进去。” 贾珍命人将门关上。 贾政先秉正看门，只见正门五间， 上面筒瓦泥鳅脊，那门栏窗？
        俱是细雕时新花样，并无朱粉涂饰。 一色水磨群墙， 下面白石台阶，凿成西番莲花样。 左右一望，雪白粉墙，下面虎皮石砌成纹理，不 落富丽俗套，自是喜欢。
        遂命开门进去。 只见一带翠嶂挡在面前。 众清客都道：“好 山，好山！” 贾政道：“非此一山，一进来园中所有之景悉入目中，更有何趣？” 众人都道：“极是。}
}
{
    \eA{According to the rites, we should, in order to obviate any shortcoming,
        request the imperial consort to deign and compose them; but if the
        honourable consort does not gaze upon the scenery with her own eyes, it will
        also be difficult for her to conceive its nature and indite upon it!}
}
{
}

\Verse{6}
{
    \zA{非胸中大有丘壑，焉能想到这里。” 说毕，往前一望，见白石 ？？ ，或如鬼怪，或似猛兽，纵横拱立。 上面苔藓斑驳，或藤萝掩映，其中微露羊 肠小径。
        贾政道：“我们就从此小径游去，回来由那一边出去，方可遍览。” 说毕，命贾珍前导，自己扶了宝玉，逶迤走进山口。 抬头忽见山上有镜面白石
        一块，正是迎面留题处。 贾政回头笑道：“诸公请看，此处题以何名方妙？”}
}
{
    \eA{And were we to wait until the arrival of her highness, to request her to
        honour the grounds with a visit, before she composes the inscriptions, such
        a wide landscape, with so many pavilions and arbours, will, without one
        character in the way of a motto, albeit it may abound with flowers, willows,
        rockeries, and streams, nevertheless in no way be able to show off its
        points of beauty to advantage."}
}
{
}

\Verse{7}
{
    \zA{众人 听说，也有说该题“叠翠”二字的，也有说该题“锦嶂”的，又有说“赛香炉”的， 又有说“小终南”的，种种名色，不止几十个。
        原来众客心中，早知贾政要试宝玉 的才情，故此只将些俗套敷衍。 宝玉也知此意。 贾政听了，便回头命宝玉拟来。 宝
        玉道：“尝听见古人说：‘编新不如述旧，刻古终胜雕今。’ 况这里并非主山正景， 原无可题，不过是探景的一进步耳。 莫如直书古人‘曲径通幽’这旧句在上，倒也
        大方。”}
}
{
    \eA{The whole party of family companions, who stood by, smiled. "Your views,
        remarkable sir," they ventured, "are excellent; but we have now a proposal
        to make. Tablets and scrolls for every locality cannot, on any account, be
        dispensed with, but they could not likewise, by any means, be determined
        upon for good!}
}
{
}

\Verse{8}
{
    \zA{众人听了，赞道：“是极，好极!二世兄天分高，才情远，不似我们读腐 了书的。” 贾政笑道：“不当过奖他。 他年小的人，不过以一知充十用，取笑罢了。
        再俟选拟。” 说着，进入石洞，只见佳木茏葱，奇花烂漫，一带清流，从花木深处泻于石隙 之下。
        再进数步，渐向北边，平坦宽豁，两边飞楼插空，雕甍绣槛，皆隐于山坳树 杪之间。 俯而视之，但见青溪泻玉，石磴穿云，白石为栏，环抱池沼，石桥三港，
        兽面衔吐。}
}
{
    \eA{Were now, for the time being, two, three or four characters fixed upon,
        harmonising with the scenery, to carry out, for form's sake, the idea, and
        were they provisionally utilised as mottoes for the lanterns, tablets and
        scrolls, and hung up, pending the arrival of her highness, and her visit
        through the grounds, when she could be requested to decide upon the devices,
        would not two exigencies be met with satisfactorily?"}
}
{
}

\Verse{9}
{
    \zA{桥上有亭，贾政与诸人到亭内坐了，问：“诸公以何题此？” 诸人都说： “当日欧阳公《醉翁亭记》有云‘有亭翼然’，就名‘翼然’罢。” 贾政笑道：“‘翼
        然’虽佳，但此亭压水而成，还须偏于水题为称。 依我拙裁，欧阳公句：‘泻于两 峰之间’，竟用他这一个‘泻’字。” 有一客道：“是极，是极。 竟是‘泻玉’二
        字妙。” 贾政拈须寻思，因叫宝玉也拟一个来。 宝玉回道：“老爷方才所说已是。}
}
{
    \eA{"Your views are perfectly correct," observed Chia Cheng, after he had heard
        their suggestion; "and we should go to-day and have a look at the place so
        as then to set to work to write the inscriptions; which, if suitable, can
        readily be used; and, if unsuitable, Yü-ts'un can then be sent for, and
        asked to compose fresh ones." The whole company smiled.}
}
{
}

\Verse{10}
{
    \zA{但如今追究了去，似乎当日欧阳公题酿泉用一‘泻’字则妥，今日此泉也用‘泻’ 字，似乎不妥。 况此处既为省亲别墅，亦当依应制之体，用此等字亦似粗陋不雅。
        求再拟蕴藉含蓄者。” 贾政笑道：“诸公听此论何如？ 方才众人编新，你说‘不如 述古’； 如今我们述古，你又说粗陋不妥。 你且说你的。”
        宝玉道：“用‘泻玉’ 二字，则不若‘沁芳’二字，岂不新雅？” 贾政拈须点头不语。 众人都忙迎合，称 赞宝玉才情不凡。 贾政道：“匾上二字容易。}
}
{
    \eA{"If you, sir, were to compose them to-day," they ventured, "they are sure to
        be excellent; and what need will there be again to wait for Yü-ts'un!" "You
        people are not aware," Chia Cheng added with a smiling countenance, "that
        I've been, even in my young days, very mediocre in the composition of
        stanzas on flowers, birds, rockeries and streams;}
}
{
}

\Verse{11}
{
    \zA{再作一副七言对来。” 宝玉四顾一望， 机上心来，乃念道： 绕堤柳借三篙翠， 隔岸花分一脉香。 贾政听了，点头微笑。 众人又称赞了一番。
        于是出亭过池，一山一石，一花一木，莫不着意观览。 忽抬头见前面一带粉垣， 数楹修舍，有千百竿翠竹遮映。 众人都道：“好个所在！” 于是大家进入，只见进
        门便是曲折游廊，阶下石子漫成甬路，上面小小三间房舍，两明一暗，里面都是合 着地步打的床几椅案。}
}
{
    \eA{and that now that I'm well up in years and have moreover the fatigue and
        trouble of my official duties, I've become in literary compositions like
        these, which require a light heart and gladsome mood, still more inapt.}
}
{
}

\Verse{12}
{
    \zA{从里间房里，又有一小门，出去却是后园，有大株梨花，阔 叶芭蕉，又有两间小小退步。 后院墙下忽开一隙，得泉一派，开沟尺许，灌入墙内，
        绕阶缘屋至前院，盘旋竹下而出。 贾政笑道：“这一处倒还好，若能月夜至此窗下 读书，也不枉虚生一世。” 说着便看宝玉，唬的宝玉忙垂了头。
        众人忙用闲话解说。 又二客说：“此处的匾该题四个字。” 贾政笑问：“那四字？” 一个道是：“淇水 遗风。” 贾政道：“俗。” 又一个道是：“睢园遗迹。”
        贾政道：“也俗。”}
}
{
    \eA{Were I even to succeed in composing any, they will unavoidably be so doltish
        and forced that they would contrariwise be instrumental in making the
        flowers, trees, garden and pavilions, through their demerits, lose in
        beauty, and present instead no pleasing feature."}
}
{
}

\Verse{13}
{
    \zA{贾珍 在旁说道：“还是宝兄弟拟一个罢。” 贾政道：“他未曾做，先要议论人家的好歹， 可见是个轻薄东西。” 众客道：“议论的是，也无奈他何。”
        贾政忙道：“休如此 纵了他。” 因说道：“今日任你狂为乱道，等说出议论来，方许你做。 方才众人说 的，可有使得的没有？”
        宝玉见问，便答道：“都似不妥。” 贾政冷笑道：“怎么 不妥？” 宝玉道：“这是第一处行幸之所，必须颂圣方可。 若用四字的匾，又有古
        人现成的，何必再做？”}
}
{
    \eA{"This wouldn't anyhow matter," remonstrated all the family companions, "for
        after perusing them we can all decide upon them together, each one of us
        recommending those he thinks best; which if excellent can be kept, and if
        faulty can be discarded; and there's nothing unfeasible about this!" "This
        proposal is most apposite," rejoined Chia Cheng. "What's more, the weather
        is, I rejoice, fine to-day;}
}
{
}

\Verse{14}
{
    \zA{贾政道：“难道‘淇水’、‘睢园’不是古人的？” 宝玉 道：“这太板了。 莫若‘有凤来仪’四字。” 众人都哄然叫妙。 贾政点头道：“畜
        生，畜生!可谓‘管窥蠡测’矣。” 因命：“再题一联来。” 宝玉便念道： 宝鼎茶闲烟尚绿， 幽窗棋罢指犹凉。 贾政摇头道：“也未见长。” 说毕，引人出来。
        方欲走时，忽想起一事来，问贾珍道：“这些院落屋宇，并几案桌椅都算有了。 还有那些帐幔帘子并陈设玩器古董，可也都是一处一处合式配就的么？”}
}
{
    \eA{so let's all go in a company and have a look." Saying this, he stood up and
        went forward, at the head of the whole party; while Chia Chen betook himself
        in advance into the garden to let every one know of their coming.}
}
{
}

\Verse{15}
{
    \zA{贾珍回 道：“那陈设的东西早已添了许多，自然临期合式陈设。 帐幔帘子，昨日听见琏兄 弟说，还不全。
        那原是一起工程之时就画了各处的图样，量准尺寸，就打发人办去 的； 想必昨日得了一半。” 贾政听了，便知此事不是贾珍的首尾，便叫人去唤贾琏。
        一时来了，贾政问他：“共有几宗？ 现今得了几宗？ 尚欠几宗？”}
}
{
    \eA{As luck would have it, Pao-yü--for he had been these last few days thinking
        of Ch'in Chung and so ceaselessly sad and wounded at heart, that dowager
        lady Chia had frequently directed the servants to take him into the new
        garden to play--made his entrance just at this very time, and suddenly
        became aware of the arrival of Chia Chen, who said to him with a smile,
        "Don't you yet run away as fast as you can?}
}
{
}

\Verse{16}
{
    \zA{贾琏见问，忙向靴 筒内取出靴掖里装的一个纸折略节来，看了一看，回道：“妆蟒洒堆、刻丝弹墨并 各色绸绫大小幔子一百二十架，昨日得了八十架，下欠四十架。
        帘子二百挂，昨日 俱得了。 外有猩猩毡帘二百挂，湘妃竹帘一百挂，金丝藤红漆竹帘一百挂，黑漆竹
        帘一百挂，五彩线络盘花帘二百挂，每样得了一半，也不过秋天都全了。 椅搭、桌 围、床裙、杌套，每分一千二百件，也有了。” 一面说，一面走，忽见青山斜阻。}
}
{
    \eA{Mr. Chia Cheng will be coming in a while." At these words, Pao-yü led off
        his nurse and the youths, and rushed at once out of the garden, like a
        streak of smoke; but as he turned a corner, he came face to face with Chia
        Cheng, who was advancing towards that direction, at the head of all the
        visitors; and as he had no time to get out of the way, the only course open
        to him was to stand on one side.}
}
{
}

\Verse{17}
{
    \zA{转过山怀中，隐隐露出一带黄泥墙，墙上皆 用稻茎掩护。 有几百枝杏花，如喷火蒸霞一般。 里面数楹茅屋，外面却是桑、榆、
        槿、柘各色树稚新条，随其曲折，编就两溜青篱。 篱外山坡之下，有一土井，旁有 桔槔辘轳之属； 下面分畦列亩，佳蔬菜花，一望无际。 贾政笑道：“倒是此处有些
        道理。 虽系人力穿凿，却入目动心，未免勾引起我归农之意。 我们且进去歇息歇息。” 说毕，方欲进去，忽见篱门外路旁有一石，亦为留题之所。}
}
{
    \eA{Chia Cheng had, of late, heard the tutor extol him by saying that he
        displayed special ability in rhyming antithetical lines, and that although
        he did not like to read his books, he nevertheless possessed some depraved
        talents, and hence it was that he was induced at this moment to promptly bid
        him follow him into the garden, with the intent of putting him to the test.}
}
{
}

\Verse{18}
{
    \zA{众人笑道：“更妙，更 妙!此处若悬匾待题，则田舍家风一洗尽矣。 立此一碣，又觉许多生色，非范石湖 田家之咏不足以尽其妙。” 贾政道：“诸公请题。”
        众人云：“方才世兄云：‘编 新不如述旧。’ 此处古人已道尽矣：莫若直书‘杏花村’为妙。” 贾政听了，笑向 贾珍道：“正亏提醒了我。
        此处都好，只是还少一个酒幌，明日竟做一个来，就依 外面村庄的式样，不必华丽，用竹竿挑在树梢头。”}
}
{
    \eA{Pao-yü could not make out what his object was, but he was compelled to
        follow. As soon as they reached the garden gate, and he caught sight of Chia
        Chen, standing on one side, along with several managers: "See that the
        garden gate is closed for a time," Chia Cheng exclaimed, "for we'll first
        see the outside and then go in."}
}
{
}

\Verse{19}
{
    \zA{贾珍答应了，又回道：“此处 竟不必养别样雀鸟，只养些鹅、鸭、鸡之类，才相称。” 贾政与众人都说好。
        贾政又向众人道：“‘杏花村’固佳，只是犯了正村名，直待请名方可。” 众 客都道：“是呀!如今虚的，却是何字样好呢？” 大家正想，宝玉却等不得了，也
        不等贾政的话，便说道：“旧诗云：‘红杏梢头挂酒旗。’ 如今莫若且题以‘杏帘 在望’四字。” 众人都道：“好个‘在望’!又暗合‘杏花村’意思。”}
}
{
    \eA{Chia Chen directed a servant to close the gate, and Chia Cheng first looked
        straight ahead of him towards the gate and espied on the same side as the
        main entrance a suite of five apartments. Above, the cylindrical tiles
        resembled the backs of mud eels. The doors, railings, windows, and frames
        were all finely carved with designs of the new fashion, and were painted
        neither in vermilion nor in white colours.}
}
{
}

\Verse{20}
{
    \zA{宝玉冷笑 道：“村名若用‘杏花’二字，便俗陋不堪了。 唐人诗里，还有‘柴门临水稻花香’， 何不用‘稻香村’的妙？” 众人听了，越发同声拍手道妙。
        贾政一声断喝：“无知 的畜生!你能知道几个古人，能记得几首旧诗，敢在老先生们跟前卖弄!方才任你胡 说，也不过试你的清浊，取笑而已，你就认真了！”
        说着，引众人步入茆堂，里面纸窗木榻，富贵气象一洗皆尽。 贾政心中自是欢 喜，却瞅宝玉道：“此处如何？” 众人见问，都忙悄悄的推宝玉教他说好。}
}
{
    \eA{The whole extent of the walls was of polished bricks of uniform colour;
        while below, the white marble on the terrace and steps was engraved with
        western foreign designs; and when he came to look to the right and to the
        left, everything was white as snow.}
}
{
}

\Verse{21}
{
    \zA{宝玉不 听人言，便应声道：“不及‘有凤来仪’多了。” 贾政听了道：“咳!无知的蠢物， 你只知朱楼画栋、恶赖富丽为佳，那里知道这清幽气象呢？
        终是不读书之过！” 宝 玉忙答道：“老爷教训的固是，但古人云‘天然’二字，不知何意？” 众人见宝玉 牛心，都怕他讨了没趣；
        今见问“天然”二字，众人忙道：“哥儿别的都明白，如 何‘天然’反要问呢？ 天然者，天之自成，不是人力之所为的。”}
}
{
    \eA{At the foot of the white-washed walls, tiger-skin pebbles were, without
        regard to pattern, promiscuously inserted in the earth in such a way as of
        their own selves to form streaks. Nothing fell in with the custom of
        gaudiness and display so much in vogue, so that he naturally felt full of
        delight;}
}
{
}

\Verse{22}
{
    \zA{宝玉道：“却又 来!此处置一田庄，分明是人力造作成的：远无邻村，近不负郭，背山无脉，临水
        无源，高无隐寺之塔，下无通市之桥，峭然孤出，似非大观，那及前数处有自然之 理、自然之趣呢？ 虽种竹引泉，亦不伤穿凿。 古人云‘天然图画’四字，正恐非其
        地而强为其地，非其山而强为其山，即百般精巧，终不相宜……”未及说完，贾政 气的喝命：“？ 出去！” 才出去，又喝命：“回来！” 命：“再题一联，若不通，
        一并打嘴巴！”}
}
{
    \eA{and, when he forthwith asked that the gate should be thrown open, all that
        met their eyes was a long stretch of verdant hills, which shut in the view
        in front of them. "What a fine hill, what a pretty hill!" exclaimed all the
        companions with one voice.}
}
{
}

\Verse{23}
{
    \zA{宝玉吓的战兢兢的，半日，只得念道： 新绿涨添浣葛处， 好云香护采芹人。 贾政听了，摇头道：“更不好。” 一面引人出来，转过山坡，穿花度柳，抚石
        依泉，过了荼？ 架，入木香棚，越牡丹亭，度芍药圃，到蔷薇院，傍芭蕉坞里盘旋 曲折。 忽闻水声潺潺，出于石洞； 上则萝薜倒垂，下则落花浮荡。
        众人都道：“好 景，好景！” 贾政道：“诸公题以何名？” 众人道：“再不必拟了，恰恰乎是‘武 陵源’三字。” 贾政笑道：“又落实了，而且陈旧。”}
}
{
    \eA{"Were it not for this one hill," Chia Cheng explained, "whatever scenery is
        contained in it would clearly strike the eye, as soon as one entered into
        the garden, and what pleasure would that have been?" "Quite so," rejoined
        all of them. "But without large hills and ravines in one's breast (liberal
        capacities), how could one attain such imagination!"}
}
{
}

\Verse{24}
{
    \zA{众人笑道：“不然就用‘秦 人旧舍’四字也罢。” 宝玉道：“越发背谬了。 ‘秦人旧舍’是避乱之意，如何使 得？ 莫若‘蓼汀花溆’四字。”
        贾政听了道：“更是胡说。” 于是贾政进了港洞，又问贾珍：“有船无船？” 贾珍道：“采莲船共四只，座 船一只，如今尚未造成。”
        贾政笑道：“可惜不得入了！” 贾珍道：“从山上盘道 也可以进去的。” 说毕，在前导引，大家攀藤抚树过去。 只见水上落花愈多，其水
        愈加清溜，溶溶荡荡，曲折萦纡。}
}
{
    \eA{After the conclusion of this remark, they cast a glance ahead of them, and
        perceived white rugged rocks looking, either like goblins, or resembling
        savage beasts, lying either crossways, or in horizontal or upright
        positions; on the surface of which grew moss and lichen with mottled hues,
        or parasitic plants, which screened off the light;}
}
{
}

\Verse{25}
{
    \zA{池边两行垂柳，杂以桃杏遮天，无一些尘土。 忽 见柳阴中又露出一个折带朱栏板桥来，度过桥去，诸路可通，便见一所清凉瓦舍， 一色水磨砖墙，清瓦花堵。
        那大主山所分之脉皆穿墙而过。 贾政道：“此处这一所 房子，无味的很。” 因而步入门时，忽迎面突出插天的大玲珑山石来，四面群绕各
        式石块，竟把里面所有房屋悉皆遮住。}
}
{
    \eA{while, slightly visible, wound, among the rocks, a narrow pathway like the
        intestines of a sheep. "If we were now to go and stroll along by this narrow
        path," Chia Cheng suggested, "and to come out from over there on our return,
        we shall have been able to see the whole grounds." Having finished speaking,
        he asked Chia Chen to lead the way; and he himself, leaning on Pao-yü,
        walked into the gorge with leisurely step.}
}
{
}

\Verse{26}
{
    \zA{且一树花木也无，只见许多异草，或有牵藤 的，或有引蔓的，或垂山岭，或穿石脚，甚至垂檐绕柱，萦砌盘阶，或如翠带飘摇，
        或如金绳蟠屈，或实若丹砂，或花如金桂，味香气馥，非凡花之可比。 贾政不禁道： “有趣!只是不大认识。” 有的说：“是薜荔藤萝。”
        贾政道：“薜荔藤萝那得有 此异香？” 宝玉道：“果然不是。 这众草中也有藤萝薜荔。 那香的是杜若蘅芜，那 一种大约是？
        兰，这一种大约是金葛，那一种是金？ 草，这一种是玉？}
}
{
    \eA{Raising his head, he suddenly beheld on the hill a block of stone, as white
        as the surface of a looking-glass, in a site which was, in very deed,
        suitable to be left for an inscription, as it was bound to meet the eye.
        "Gentlemen," Chia Cheng observed, as he turned his head round and smiled,
        "please look at this spot. What name will it be fit to give it?"}
}
{
}

\Verse{27}
{
    \zA{藤，红的自 然是紫芸，绿的定是青芷。 想来那《离骚》、《文选》所有的那些异草：有叫作什 么霍纳姜汇的，也有叫作什么纶组紫绛的。
        还有什么石帆、清松、扶留等样的，见 于左太冲《吴都赋》。 又有叫作什么绿荑的，还有什么丹椒、蘼芜、风莲，见于《蜀 都赋》。
        如今年深岁改，人不能识，故皆象形夺名，渐渐的唤差了，也是有的。” 未及说完，贾政喝道：“谁问你来？” 唬的宝玉倒退，不敢再说。}
}
{
    \eA{When the company heard his remark, some maintained that the two words
        "Heaped verdure" should be written; and others upheld that the device should
        be "Embroidered Hill." Others again suggested: "Vying with the Hsiang Lu;"
        and others recommended "the small Chung Nan." And various kinds of names
        were proposed, which did not fall short of several tens.}
}
{
}

\Verse{28}
{
    \zA{贾政因见两边俱是超手游廊，便顺着游廊步入，只见上面五间清厦，连着卷棚， 四面出廊，绿窗油壁，更比前清雅不同。 贾政叹道：“此轩中煮茗操琴，也不必再
        焚香了。 此造却出意外，诸公必有佳作新题以颜其额，方不负此。” 众人笑道：“莫 若‘兰风蕙露’贴切了。” 贾政道：“也只好用这四字。 其联云何？”
        一人道：“我 想了一对，大家批削改正。 道是：‘麝兰芳霭斜阳院，杜若香飘明月洲。’” 众人 道：“妙则妙矣!只是‘斜阳’二字不妥。”}
}
{
    \eA{All the visitors had been, it must be explained, aware at an early period of
        the fact that Chia Cheng meant to put Pao-yü's ability to the test, and for
        this reason they merely proposed a few combinations in common use. But of
        this intention, Pao-yü himself was likewise cognizant. After listening to
        the suggestions, Chia Cheng forthwith turned his head round and bade Pao-yü
        think of some motto.}
}
{
}

\Verse{29}
{
    \zA{那人引古诗“蘼芜满院泣斜阳”句， 众人云：“颓丧，颓丧！” 又一人道：“我也有一联，诸公评阅评阅。” 念道：“三 径香风飘玉蕙，一庭明月照金兰。”
        贾政拈须沉吟，意欲也题一联。 忽抬头见宝玉 在旁不敢作声，因喝道：“怎么你应说话时又不说了!还要等人请教你不成？” 宝
        玉听了回道：“此处并没有什么‘兰麝’、‘明月’、‘洲渚’之类，若要这样着 迹说来，就题二百联也不能完。” 贾政道：“谁按着你的头，教你必定说这些字样
        呢？”}
}
{
    \eA{"I've often heard," Pao-yü replied, "that writers of old opine that it's
        better to quote an old saying than to compose a new one; and that an old
        engraving excels in every respect an engraving of the present day.}
}
{
}

\Verse{30}
{
    \zA{宝玉道：“如此说，则匾上莫若‘蘅芷清芬’四字。 对联则是：‘吟成豆蔻 诗犹艳，睡足荼？ 梦亦香。’” 贾政笑道：“这是套的‘书成蕉叶文犹绿’，不足
        为奇。” 众人道：“李太白‘凤凰台’之作，全套‘黄鹤楼’。 只要套得妙。 如今 细评起来，方才这一联竟比‘书成蕉叶’尤觉幽雅活动。”
        贾政笑道：“岂有此理。” 说着，大家出来。 走不多远，则见崇阁巍峨，层楼高起，面面琳宫合抱，迢迢 复道萦纡。 青松拂檐，玉兰绕砌； 金辉兽面，彩焕螭头。}
}
{
    \eA{What's more, this place doesn't constitute the main hill or the chief
        feature of the scenery, and is really no site where any inscription should
        be put, as it no more than constitutes the first step in the inspection of
        the landscape. Won't it be well to employ the exact text of an old writer
        consisting of 'a tortuous path leading to a secluded (nook).'}
}
{
}

\Verse{31}
{
    \zA{贾政道：“这是正殿了。 只是太富丽了些！” 众人都道：“要如此方是。 虽然贵妃崇尚节俭，然今日之尊， 礼仪如此，不为过也。”
        一面说，一面走，只见正面现出一座玉石牌坊，上面龙蟠 螭护，玲珑凿就。 贾政道：“此处书以何文？” 众人道：“必是‘蓬莱仙境’方妙。” 贾政摇头不语。
        宝玉见了这个所在，心中忽有所动，寻思起来，倒像在那里见过的 一般，却一时想不起那年那日的事了。 贾政又命他题咏，宝玉只顾细思前景，全无 心于此了。}
}
{
    \eA{This line of past days would, if inscribed, be, in fact, liberal to boot."
        After listening to the proposed line, they all sang its praise. "First-rate!
        excellent!" they cried, "the natural talents of your second son, dear
        friend, are lofty; his mental capacity is astute; he is unlike ourselves,
        who have read books but are simple fools." "You shouldn't," urged Chia Cheng
        smilingly, "heap upon him excessive praise;}
}
{
}

\Verse{32}
{
    \zA{众人不知其意，只当他受了这半日折磨，精神耗散，才尽词穷了，再要 牛难逼迫着了急，或生出事来，倒不便。 遂忙都劝贾政道：“罢了，明日再题罢了。”
        贾政心中也怕贾母不放心，遂冷笑道：“你这畜生，也竟有不能之时了。 也罢，限 你一日，明日题不来，定不饶你。 这是第一要紧处所，要好生作来！”
        说着，引人出来，再一观望，原来自进门至此，才游了十之五六。 又值人来回， 有雨村处遣人回话。 贾政笑道：“此数处不能游了。}
}
{
    \eA{he's young in years, and merely knows one thing which he turns to the use of
        ten purposes; you should laugh at him, that's all; but we can by and by
        choose some device."}
}
{
}

\Verse{33}
{
    \zA{虽如此，到底从那一边出去， 也可略观大概。” 说着，引客行来，至一大桥，水如晶帘一般奔入。 原来这桥边是 通外河之闸，引泉而入者。
        贾政因问：“此闸何名？” 宝玉道：“此乃沁芳源之正 流，即名‘沁芳闸’。” 贾政道：“胡说，偏不用‘沁芳’二字。”
        于是一路行来，或清堂，或茅舍，或堆石为垣，或编花为门，或山下得幽尼佛 寺，或林中藏女道丹房，或长廊曲洞，或方厦圆亭：贾政皆不及进去。}
}
{
    \eA{As he spoke, he entered the cave, where he perceived beautiful trees with
        thick foliage, quaint flowers in lustrous bloom, while a line of limpid
        stream emanated out of a deep recess among the flowers and trees, and oozed
        down through the crevice of the rock.}
}
{
}

\Verse{34}
{
    \zA{因半日未尝 歇息，腿酸脚软，忽又见前面露出一所院落来，贾政道：“到此可要歇息歇息了。”
        说着一径引入，绕着碧桃花，穿过竹篱花障编就的月洞门，俄见粉垣环护，绿柳周 垂。 贾政与众人进了门，两边尽是游廊相接，院中点衬几块山石，一边种几本芭蕉，
        那一边是一树西府海棠，其势若伞，丝垂金缕，葩吐丹砂。 众人都道：“好花，好 花!海棠也有，从没见过这样好的。”}
}
{
    \eA{Progressing several steps further in, they gradually faced the northern
        side, where a stretch of level ground extended far and wide, on each side of
        which soared lofty buildings, intruding themselves into the skies, whose
        carved rafters and engraved balustrades nestled entirely among the
        depressions of the hills and the tops of the trees.}
}
{
}

\Verse{35}
{
    \zA{贾政道：“这叫做‘女儿棠’，乃是外国之 种，俗传出‘女儿国’，故花最繁盛，――亦荒唐不经之说耳。” 众人道：“毕竟 此花不同，‘女国’之说，想亦有之。”
        宝玉云：“大约骚人咏士以此花红若施脂， 弱如扶病，近乎闺阁风度，故以‘女儿’命名，世人以讹传讹，都未免认真了。” 众人都说：“领教!妙解！”
        一面说话，一面都在廊下榻上坐了。 贾政因道：“想 几个什么新鲜字来题？” 一客道：“‘蕉鹤’二字妙。” 又一个道：“‘崇光泛彩’ 方妙。”}
}
{
    \eA{They lowered their eyes and looked, and beheld a pure stream flowing like
        jade, stone steps traversing the clouds, a balustrade of white marble
        encircling the pond in its embrace, and a stone bridge with three archways,
        the animals upon which had faces disgorging water from their mouths. A
        pavilion stood on the bridge, and in this pavilion Chia Chen and the whole
        party went and sat.}
}
{
}

\Verse{36}
{
    \zA{贾政与众人都道：“好个‘崇光泛彩’！” 宝玉也道：“妙。” 又说：“只 是可惜了！” 众人问：“如何可惜？” 宝玉道：“此处蕉棠两植，其意暗蓄‘红’
        ‘绿’二字在内，若说一样，遗漏一样，便不足取。” 贾政道：“依你如何？” 宝 玉道：“依我，题‘红香绿玉’四字，方两全其美。”
        贾政摇头道：“不好，不好！” 说着，引人进入房内。 只见其中收拾的与别处不同，竟分不出间隔来。}
}
{
    \eA{"Gentlemen," he inquired, "what shall we write about this?" "In the record,"
        they all replied, "of the 'Drunken Old Man's Pavilion,' written in days of
        old by Ou Yang, appears this line: 'There is a pavilion pinioned-like,' so
        let us call this 'the pinioned-like pavilion,' and finish." "Pinioned-like,"
        observed Chia Cheng smiling, "is indeed excellent;}
}
{
}

\Verse{37}
{
    \zA{原来四 面皆是雕空玲珑木板，或“流云百蝠”，或“岁寒三友”，或山水人物，或翎毛花
        卉，或集锦，或博古，或万福万寿，各种花样，皆是名手雕镂五彩，销金嵌玉的。 一？ 一？ ，或贮书，或设鼎，或安置笔砚，或供设瓶花，或安放盆景。 其？
        式样或 圆或方，或葵花蕉叶，或连环半璧，真是花团锦簇，剔透玲珑。 倏尔五色纱糊，竟 系小窗； 倏尔彩绫轻覆，竟系幽户。
        且满墙皆是随依古董玩器之形抠成的槽子，如 琴、剑、悬瓶之类，俱悬于壁，却都是与壁相平的。}
}
{
    \eA{but this pavilion is constructed over the water, and there should, after
        all, be some allusion to the water in the designation. My humble opinion is
        that of the line in Ou Yang's work, '(the water) drips from between the two
        peaks,' we should only make use of that single word 'drips.'" "First-rate!"
        rejoined one of the visitors, "capital!}
}
{
}

\Verse{38}
{
    \zA{众人都赞：“好精致!难为怎 么做的！” 原来贾政走进来了，未到两层，便都迷了旧路，左瞧也有门可通，右瞧
        也有窗隔断，及到跟前，又被一架书挡住，回头又有窗纱明透门径。 及至门前，忽 见迎面也进来了一起人，与自己的形相一样，――却是一架大玻璃镜。 转过镜去，
        一发见门多了。 贾珍笑道：“老爷随我来，从这里出去就是后院，出了后院倒比先 近了。” 引着贾政及众人转了两层纱厨，果得一门出去，院中满架蔷薇。
        转过花障， 只见青溪前阻。}
}
{
    \eA{but what would really be appropriate are the two characters 'dripping
        jadelike.'" Chia Chen pulled at his moustache, as he gave way to reflection;
        after which, he asked Pao-yü to also propose one himself. "What you, sir,
        suggested a while back," replied Pao-yü, "will do very well;}
}
{
}

\Verse{39}
{
    \zA{众人诧异：“这水又从何而来？” 贾珍遥指道：“原从那闸起流至 那洞口，从东北山凹里引到那村庄里，又开一道岔口，引至西南上，共总流到这里，
        仍旧合在一处，从那墙下出去。” 众人听了，都道：“神妙之极！” 说着，忽见大 山阻路，众人都迷了路，贾珍笑道：“跟我来。” 乃在前导引，众人随着，由山脚
        下一转，便是平坦大路，豁然大门现于面前，众人都道：“有趣，有趣!搜神夺巧， 至于此极！” 于是大家出来。}
}
{
    \eA{but if we were now to sift the matter thoroughly, the use of the single word
        'drip' by Ou Yang, in his composition about the Niang spring, would appear
        quite apposite; while the application, also on this occasion, to this
        spring, of the character 'drip' would be found not quite suitable.}
}
{
}

\Verse{40}
{
    \zA{那宝玉一心只记挂着里边姊妹们，又不见贾政吩咐，只得跟到书房。 贾政忽想 起来道：“你还不去，看老太太惦记你。 难道还逛不足么？” 宝玉方退了出来。 至
        院外，就有跟贾政的小厮上来抱住，说道：“今日亏了老爷喜欢，方才老太太打发 人出来问了几遍，我们回说老爷喜欢； 要不然，老太太叫你进去了，就不得展才了。
        人人都说你才那些诗比众人都强，今儿得了彩头，该赏我们了。” 宝玉笑道：“每 人一吊。” 众人道：“谁没见那一吊钱!把这荷包赏了罢。”}
}
{
    \eA{Moreover, seeing that this place is intended as a separate residence (for
        the imperial consort), on her visit to her parents, it is likewise
        imperative that we should comply with all the principles of etiquette, so
        that were words of this kind to be used, they would besides be coarse and
        inappropriate; and may it please you to fix upon something else more
        recondite and abstruse."}
}
{
}

\Verse{41}
{
    \zA{说着，一个个都上来 解荷包，解扇袋，不容分说，将宝玉所佩之物，尽行解去。 又道：“好生送上去罢。” 一个个围绕着，送至贾母门前。
        那时贾母正等着他，见他来了，知道不曾难为他， 心中自是喜欢。 少时袭人倒了茶来，见身边佩物一件不存，因笑道：“带的东西必又是那起没
        脸的东西们解了去了。” 黛玉听说，走过来一瞧，果然一件没有，因向宝玉道：“我 给你的那个荷包也给他们了？ 你明儿再想我的东西，可不能够了！”}
}
{
    \eA{"What do you, gentlemen, think of this argument?" Chia Cheng remarked
        sneeringly. "A little while ago, when the whole company devised something
        original, you observed that it would be better to quote an old device; and
        now that we have quoted an old motto, you again maintain that it's coarse
        and inappropriate! But you had better give us one of yours."}
}
{
}

\Verse{42}
{
    \zA{说毕，生气回 房，将前日宝玉嘱咐他没做完的香袋儿，拿起剪子来就铰。 宝玉见他生气，便忙赶 过来，早已剪破了。
        宝玉曾见过这香袋，虽未完工，却十分精巧，无故剪了，却也 可气。 因忙把衣领解了，从里面衣襟上将所系荷包解下来了，递与黛玉道：“你瞧 瞧，这是什么东西？
        我何从把你的东西给人来着？” 黛玉见他如此珍重，带在里面， 可知是怕人拿去之意，因此自悔莽撞剪了香袋，低着头一言不发。 宝玉道：“你也
        不用铰，我知你是懒怠给我东西。}
}
{
    \eA{"If two characters like 'dripping jadelike' are to be used," Pao-yü
        explained, "it would be better then to employ the two words 'Penetrating
        Fragrance,' which would be unique and excellent, wouldn't they?" Chia Cheng
        pulled his moustache, nodded his head and did not utter a word; whereupon
        the whole party hastily pressed forward with one voice to eulogize Pao-yü's
        acquirements as extraordinary.}
}
{
}

\Verse{43}
{
    \zA{我连这荷包奉还，何如？” 说着掷向他怀中而去。 黛玉越发气的哭了，拿起荷包又铰。 宝玉忙回身抢住，笑道：“好妹妹饶了他罢！”
        黛玉将剪子一摔，拭泪说道：“你不用合我好一阵歹一阵的，要恼就撂开手。” 说 着赌气上床，面向里倒下拭泪。 禁不住宝玉上来“妹妹”长“妹妹”短赔不是。
        前面贾母一片声找宝玉。 众人回说：“在林姑娘房里。” 贾母听说道：“好， 好!让他姐妹们一处玩玩儿罢。 才他老子拘了他这半天，让他松泛一会子罢。}
}
{
    \eA{"The selection of two characters for the tablet is an easy matter,"
        suggested Chia Cheng, "but now go on and compose a pair of antithetical
        phrases with seven words in each." Pao-yü cast a glance round the four
        quarters, when an idea came into his head, and he went on to recite:  The
        willows, which enclose the shore, the green borrow from three  bamboos;}
}
{
}

\Verse{44}
{
    \zA{只别 叫他们拌嘴。” 众人答应着。 黛玉被宝玉缠不过，只得起来道：“你的意思不叫我安生，我就离了你。” 说 着往外就走。
        宝玉笑道：“你到那里我跟到那里。” 一面仍拿着荷包来带上。 黛玉 伸手抢道：“你说不要，这会子又带上，我也替你怪臊的！” 说着“嗤”的一声笑 了。
        宝玉道：“好妹妹，明儿另替我做个香袋儿罢！” 黛玉道：“那也瞧我的高兴 罢了。” 一面说，一面二人出房，到王夫人上房中去了。 可巧宝钗也在那里。
        此时王夫人那边热闹非常。}
}
{
    \eA{On banks apart, the flowers asunder grow, yet one perfume they give. Upon
        hearing these lines, Chia Cheng gave a faint smile, as he nodded his head,
        whilst the whole party went on again to be effusive in their praise. But
        forthwith they issued from the pavilions, and crossed the pond,
        contemplating with close attention each elevation, each stone, each flower,
        or each tree.}
}
{
}

\Verse{45}
{
    \zA{原来贾蔷已从姑苏采买了十二个女孩子、并聘了教 习以及行头等事来了，那时薛姨妈另于东北上一所幽静房舍居住，将梨香院另行修 理了，就令教习在此教演女戏；
        又另派了家中旧曾学过歌唱的众女人们，―――如 今皆是皤然老妪，着他们带领管理。 其日月出入银钱等事，以及诸凡大小所需之物 料帐目，就令贾蔷总理。
        又有林之孝来回：“采访聘买得十二个小尼姑、小道姑，都到了。 连新做的二 十分道袍也有了。}
}
{
    \eA{And as suddenly they raised their heads, they caught sight, in front of
        them, of a line of white wall, of numbers of columns, and beautiful
        cottages, where flourished hundreds and thousands of verdant bamboos, which
        screened off the rays of the sun. "What a lovely place!" they one and all
        exclaimed.}
}
{
}

\Verse{46}
{
    \zA{外又有一个带发修行的，本是苏州人氏，祖上也是读书仕宦之家， 因自幼多病，买了许多替身，皆不中用，到底这姑娘入了空门，方才好了，所以带 发修行。
        今年十八岁，取名妙玉。 如今父母俱已亡故，身边只有两个老嬷嬷、一个 小丫头伏侍，文墨也极通，经典也极熟，模样又极好。 因听说长安都中有观音遗迹
        并贝叶遗文，去年随了师父上来，现在西门外牟尼院住着。 他师父精演先天神数， 于去冬圆寂了。 遗言说他：‘不宜回乡，在此静候，自有结果。’}
}
{
    \eA{Speedily the whole company penetrated inside, perceiving, as soon as they
        had entered the gate, a zigzag arcade, below the steps of which was a raised
        pathway, laid promiscuously with stones, and on the furthest part stood a
        diminutive cottage with three rooms, two with doors leading into them and
        one without.}
}
{
}

\Verse{47}
{
    \zA{所以未曾扶灵回 去。” 王夫人便道：“这样我们何不接了他来？” 林之孝家的回道：“若请他，他 说：‘侯门公府，必以贵势压人，我再不去的。’”
        王夫人道：“他既是宦家小姐， 自然要性傲些。 就下个请帖请他何妨。” 林之孝家的答应着出去，叫书启相公写个 请帖去请妙玉，次日遣人备车轿去接。
        不知后来如何，且听下回分解。 (本章完)}
}
{
    \eA{Everything in the interior, in the shape of beds, teapoys, chairs and
        tables, were made to harmonise with the space available. Leading out of the
        inner room of the cottage was a small door from which, as they egressed,
        they found a back-court with lofty pear trees in blossom and banana trees,
        as well as two very small retiring back-courts.}
}
{
}

\Verse{48}
{
    \zA{}
}
{
    \eA{At the foot of the wall, unexpectedly became visible an aperture where was a
        spring, for which a channel had been opened scarcely a foot or so wide, to
        enable it to run inside the wall. Winding round the steps, it skirted the
        buildings until it reached the front court, where it coiled and curved,
        flowing out under the bamboos. "This spot," observed Chia Cheng full of
        smiles, "is indeed pleasant! and could one, on a moonlight night, sit under
        the window and study, one would not spend a whole lifetime in vain!" As he
        said this, he quickly cast a glance at Pao-yü, and so terrified did Pao-yü
        feel that he hastily drooped his head.}
}
{
}

\Verse{49}
{
    \zA{}
}
{
    \eA{The whole company lost no time in choosing some irrelevant talk to turn the
        conversation, and two of the visitors prosecuted their remarks by adding
        that on the tablet, in this spot, four characters should be inscribed.
        "Which four characters?" Chia Cheng inquired, laughingly. "The bequeathed
        aspect of the river Ch'i!" suggested one of them. "It's commonplace,"
        observed Chia Cheng. Another person recommended "the remaining vestiges of
        the Chü Garden." "This too is commonplace!" replied Chia Cheng. "Let brother
        Pao-yü again propound one!" interposed Chia Chen, who stood by. "Before he
        composes any himself," Chia Cheng continued, "his wont is to first discuss
        the pros and cons of those of others; so it's evident that he's an impudent
        fellow!"}
}
{
}

\Verse{50}
{
    \zA{}
}
{
    \eA{"He's most reasonable in his arguments," all the visitors protested, "and
        why should he be called to task?" "Don't humour him so much!" Chia Cheng
        expostulated. "I'll put up for to-day," he however felt constrained to tell
        Pao-yü, "with your haughty manner, and your rubbishy speech, so that after
        you have, to begin with, given us your opinion, you may next compose a
        device. But tell me, are there any that will do among the mottoes suggested
        just now by all the gentlemen?" "They all seem to me unsuitable!" Pao-yü did
        not hesitate to say by way of reply to this question. Chia Cheng gave a
        sardonic smile. "How all unsuitable?" he exclaimed.}
}
{
}

\Verse{51}
{
    \zA{}
}
{
    \eA{"This," continued Pao-yü, "is the first spot which her highness will honour
        on her way, and there should be inscribed, so that it should be appropriate,
        something commending her sacred majesty. But if a tablet with four
        characters has to be used, there are likewise devices ready at hand, written
        by poets of old; and what need is there to compose any more?" "Are forsooth
        the devices 'the river Ch'i and the Chu Garden' not those of old authors?"
        insinuated Chia Cheng. "They are too stiff," replied Pao-yü. "Would not the
        four characters: 'a phoenix comes with dignified air,' be better?" With
        clamorous unanimity the whole party shouted: "Excellent:" and Chia Cheng
        nodding his head; "You beast, you beast!"}
}
{
}

\Verse{52}
{
    \zA{}
}
{
    \eA{he ejaculated, "it may well be said about you that you see through a thin
        tube and have no more judgment than an insect! Compose another stanza," he
        consequently bade him; and Pao-yü recited:  In the precious tripod kettle,
        tea is brewed, but green is still the  smoke! O'er is the game of chess by
        the still window, but the fingers are yet  cold. Chia Cheng shook his head.
        "Neither does this seem to me good!" he said; and having concluded this
        remark he was leading the company out, when just as he was about to proceed,
        he suddenly bethought himself of something. "The several courts and
        buildings and the teapoys, sideboards, tables and chairs," he added, "may be
        said to be provided for. But there are still all those curtains, screens and
        portieres, as well as the furniture, nicknacks and curios;}
}
{
}

\Verse{53}
{
    \zA{}
}
{
    \eA{and have they too all been matched to suit the requirements of each place?"
        "Of the things that have to be placed about," Chia Chen explained, a good
        number have, at an early period, been added, and of course when the time
        comes everything will be suitably arranged. As for the curtains, screens,
        and portieres, which have to be hung up, I heard yesterday brother Lien say
        that they are not as yet complete, that when the works were first taken in
        hand, the plan of each place was drawn, the measurements accurately
        calculated and some one despatched to attend to the things, and that he
        thought that yesterday half of them were bound to come in.}
}
{
}

\Verse{54}
{
    \zA{}
}
{
    \eA{Chia Cheng, upon hearing this explanation, readily remembered that with all
        these concerns Chia Chen had nothing to do; so that he speedily sent some
        one to go and call Chia Lien. Having arrived in a short while, "How many
        sorts of things are there in all?" Chia Cheng inquired of him. "Of these how
        many kinds have by this time been got ready? and how many more are short?"
        At this question, Chia Lien hastily produced, from the flaps of his boot, a
        paper pocket-book, containing a list, which he kept inside the tops of his
        boot. After perusing it and reperusing it, he made suitable reply.}
}
{
}

\Verse{55}
{
    \zA{}
}
{
    \eA{"Of the hundred and twenty curtains," he proceeded, "of stiff spotted silks,
        embroidered with dragons in relief, and of the curtains large and small, of
        every kind of damask silk, eighty were got yesterday, so that there still
        remain forty of them to come. The two portieres were both received
        yesterday; and besides these, there are the two hundred red woollen
        portieres, two hundred portieres of Hsiang Fei bamboo; two hundred
        door-screens of rattan, with gold streaks, and of red lacquered bamboo; two
        hundred portieres of black lacquered rattan; two hundred door-screens of
        variegated thread-netting with clusters of flowers. Of each of these kinds,
        half have come in, but the whole lot of them will be complete no later than
        autumn.}
}
{
}

\Verse{56}
{
    \zA{}
}
{
    \eA{Antimacassars, table-cloths, flounces for the beds, and cushions for the
        stools, there are a thousand two hundred of each, but these likewise are
        ready and at hand." As he spoke, they proceeded outwards, but suddenly they
        perceived a hill extending obliquely in such a way as to intercept the
        passage; and as they wound round the curve of the hill faintly came to view
        a line of yellow mud walls, the whole length of which was covered with paddy
        stalks for the sake of protection, and there were several hundreds of
        apricot trees in bloom, which presented the appearance of being fire,
        spurted from the mouth, or russet clouds, rising in the air. Inside this
        enclosure, stood several thatched cottages.}
}
{
}

\Verse{57}
{
    \zA{}
}
{
    \eA{Outside grew, on the other hand, mulberry trees, elms, mallows, and silkworm
        oaks, whose tender shoots and new twigs, of every hue, were allowed to bend
        and to intertwine in such a way as to form two rows of green fence. Beyond
        this fence and below the white mound, was a well, by the side of which stood
        a well-sweep, windlass and such like articles; the ground further down being
        divided into parcels, and apportioned into fields, which, with the fine
        vegetables and cabbages in flower, presented, at the first glance, the
        aspect of being illimitable. "This is," Chia Cheng observed chuckling, "the
        place really imbued with a certain amount of the right principle;}
}
{
}

\Verse{58}
{
    \zA{}
}
{
    \eA{and laid out, though it has been by human labour, yet when it strikes my
        eye, it so moves my heart, that it cannot help arousing in me the wish to
        return to my native place and become a farmer. But let us enter and rest a
        while." As he concluded these words, they were on the point of walking in,
        when they unexpectedly discerned a stone, outside the trellis gate, by the
        roadside, which had also been left as a place on which to inscribe a motto.
        "Were a tablet," argued the whole company smilingly, "put up high in a spot
        like this, to be filled up by and by, the rustic aspect of a farm would in
        that case be completely done away with;}
}
{
}

\Verse{59}
{
    \zA{}
}
{
    \eA{and it will be better, yea far better to erect this slab on the ground, as
        it will further make manifest many points of beauty. But unless a motto
        could be composed of the same excellence as that in Fan Shih-hu's song on
        farms, it will not be adequate to express its charms!" "Gentlemen," observed
        Chia Cheng, "please suggest something." "A short while back," replied the
        whole company, "your son, venerable brother, remarked that devising a new
        motto was not equal to quoting an old one, and as sites of this kind have
        been already exhausted by writers of days of old, wouldn't it be as well
        that we should straightway call it the 'apricot blossom village?' and this
        will do splendidly."}
}
{
}

\Verse{60}
{
    \zA{}
}
{
    \eA{When Chia Cheng heard this remark, he smiled and said, addressing himself to
        Chia Chen: "This just reminds me that although this place is perfect in
        every respect, there's still one thing wanting in the shape of a wine board;
        and you had better then have one made to-morrow on the very same pattern as
        those used outside in villages; and it needn't be anything gaudy, but hung
        above the top of a tree by means of bamboos." Chia Chen assented. "There's
        no necessity," he went on to explain, "to keep any other birds in here, but
        only to rear a few geese, ducks, fowls and such like; as in that case they
        will be in perfect keeping with the place." "A splendid idea!" Chia Cheng
        rejoined, along with all the party.}
}
{
}

\Verse{61}
{
    \zA{}
}
{
    \eA{"'Apricot blossom village' is really first-rate," continued Chia Cheng as he
        again addressed himself to the company; "but the only thing is that it
        encroaches on the real designation of the village; and it will be as well to
        wait (until her highness comes), when we can request her to give it a name."
        "Certainly!" answered the visitors with one voice; "but now as far as a name
        goes, for mere form, let us all consider what expressions will be suitable
        to employ." Pao-yü did not however give them time to think;}
}
{
}

\Verse{62}
{
    \zA{}
}
{
    \eA{nor did he wait for Chia Cheng's permission, but suggested there and then:
        "In old poetical works there's this passage: 'At the top of the red apricot
        tree hangs the flag of an inn,' and wouldn't it be advisable, on this
        occasion, to temporarily adopt the four words: 'the sign on the apricot tree
        is visible'?" "'Is visible' is excellent," suggested the whole number of
        them, "and what's more it secretly accords with the meaning implied by
        'apricot blossom village.'" "Were the two words 'apricot blossom' used for
        the name of the village, they would be too commonplace and unsuitable;"}
}
{
}

\Verse{63}
{
    \zA{}
}
{
    \eA{added Pao-yü with a sardonic grin, "but there's another passage in the works
        of a poet of the T'ang era: 'By the wooden gate near the water the
        corn-flower emits its fragrance;' and why not make use of the motto 'corn
        fragrance village,' which will be excellent?" When the company heard his
        proposal, they, with still greater vigour, unanimously combined in crying
        out "Capital!" as they clapped their hands. Chia Cheng, with one shout,
        interrupted their cries, "You ignorant child of wrath!" he ejaculated; "how
        many old writers can you know, and how many stanzas of ancient poetical
        works can you remember, that you will have the boldness to show off in the
        presence of all these experienced gentlemen?}
}
{
}

\Verse{64}
{
    \zA{}
}
{
    \eA{(In allowing you to give vent to) all the nonsense you uttered my object was
        no other than to see whether your brain was clear or muddled; and all for
        fun's sake, that's all; and lo, you've taken things in real earnest!" Saying
        this, he led the company into the interior of the hall with the mallows. The
        windows were pasted with paper, and the bedsteads made of wood, and all
        appearance of finery had been expunged, and Chia Cheng's heart was naturally
        much gratified; but nevertheless, scowling angrily at Pao-yü, "What do you
        think of this place?" he asked. When the party heard this question, they all
        hastened to stealthily give a nudge to Pao-yü, with the express purpose of
        inducing him to say it was nice;}
}
{
}

\Verse{65}
{
    \zA{}
}
{
    \eA{but Pao-yü gave no ear to what they all urged. "It's by far below the spot,"
        he readily replied, "designated 'a phoenix comes with dignified air.'" "You
        ignorant stupid thing!" exclaimed Chia Cheng at these words; "what you
        simply fancy as exquisite, with that despicable reliance of yours upon
        luxury and display, are two-storied buildings and painted pillars! But how
        can you know anything about this aspect so pure and unobtrusive, and this is
        all because of that failing of not studying your books!" "Sir," hastily
        answered Pao-yü, "your injunctions are certainly correct; but men of old
        have often made allusion to 'natural;' and what is, I wonder, the import of
        these two characters?"}
}
{
}

\Verse{66}
{
    \zA{}
}
{
    \eA{The company had perceived what a perverse mind Pao yü possessed, and they
        one and all were much surprised that he should be so silly beyond the
        possibility of any change; and when now they heard the question he asked,
        about the two characters representing "natural," they, with one accord,
        speedily remarked, "Everything else you understand, and how is it that on
        the contrary you don't know what 'natural' implies? The word 'natural' means
        effected by heaven itself and not made by human labour." "Well, just so,"
        rejoined Pao-yü; "but the farm, which is laid out in this locality, is
        distinctly the handiwork of human labour; in the distance, there are no
        neighbouring hamlets; near it, adjoin no wastes; though it bears a hill, the
        hill is destitute of streaks; though it be close to water, this water has no
        spring;}
}
{
}

\Verse{67}
{
    \zA{}
}
{
    \eA{above, there is no pagoda nestling in a temple; below, there is no bridge
        leading to a market; it rises abrupt and solitary, and presents no grand
        sight! The palm would seem to be carried by the former spot, which is imbued
        with the natural principle, and possesses the charms of nature; for, though
        bamboos have been planted in it, and streams introduced, they nevertheless
        do no violence to the works executed. 'A natural landscape,' says, an
        ancient author in four words; and why? Simply because he apprehended that
        what was not land, would, by forcible ways, be converted into land; and that
        what was no hill would, by unnatural means, be raised into a hill. And
        ingenious though these works might be in a hundred and one ways, they
        cannot, after all, be in harmony."...}
}
{
}

\Verse{68}
{
    \zA{}
}
{
    \eA{But he had no time to conclude, as Chia Cheng flew into a rage. "Drive him
        off," he shouted; (but as Pao-yü) was on the point of going out, he again
        cried out: "Come back! make up," he added, "another couplet, and if it isn't
        clear, I'll for all this give you a slap on your mouth." Pao-yü had no
        alternative but to recite as follows:  A spot in which the "Ko" fibre to
        bleach, as the fresh tide doth swell  the waters green! A beauteous halo and
        a fragrant smell the man encompass who the cress  did pluck! Chia Cheng,
        after this recital, nodded his head. "This is still worse!"}
}
{
}

\Verse{69}
{
    \zA{}
}
{
    \eA{he remarked, but as he reproved him, he led the company outside, and winding
        past the mound, they penetrated among flowers, and wending their steps by
        the willows, they touched the rocks and lingered by the stream. Passing
        under the trellis with yellow roses, they went into the shed with white
        roses; they crossed by the pavilion with peonies, and walked through the
        garden, where the white peony grew; and entering the court with the cinnamon
        roses, they reached the island of bananas. As they meandered and zigzagged,
        suddenly they heard the rustling sound of the water, as it came out from a
        stone cave, from the top of which grew parasitic plants drooping downwards,
        while at its bottom floated the fallen flowers. "What a fine sight!" they
        all exclaimed; "what beautiful scenery!"}
}
{
}

\Verse{70}
{
    \zA{}
}
{
    \eA{"Gentlemen," observed Chia Cheng, "what name do you propose for this place?"
        "There's no further need for deliberation," the company rejoined; "for this
        is just the very spot fit for the three words 'Wu Ling Spring.'" "This too
        is matter-of-fact!" Chia Cheng objected laughingly, "and likewise
        antiquated." "If that won't do," the party smiled, "well then what about the
        four characters implying 'An old cottage of a man of the Ch'in dynasty?'"
        "This is still more exceedingly plain!" interposed Pao-yü. "'The old cottage
        of a man of the Ch'in dynasty' is meant to imply a retreat from revolution,
        and how will it suit this place?}
}
{
}

\Verse{71}
{
    \zA{}
}
{
    \eA{Wouldn't the four characters be better denoting 'an isthmus with smart weed,
        and a stream with flowers'?" When Chia Cheng heard these words, he
        exclaimed: "You're talking still more stuff and nonsense?" and forthwith
        entering the grotto, Chia Cheng went on to ask of Chia Chen, "Are there any
        boats or not?" "There are to be," replied Chia Chen, "four boats in all from
        which to pick the lotus, and one boat for sitting in; but they haven't now
        as yet been completed." "What a pity!" Chia Cheng answered smilingly, "that
        we cannot go in." "But we could also get into it by the tortuous path up the
        hill," Chia Chen ventured;}
}
{
}

\Verse{72}
{
    \zA{}
}
{
    \eA{and after finishing this remark, he walked ahead to show the way, and the
        whole party went over, holding on to the creepers, and supporting themselves
        by the trees, when they saw a still larger quantity of fallen leaves on the
        surface of the water, and the stream itself, still more limpid, gently and
        idly meandering along on its circuitous course. By the bank of the pond were
        two rows of weeping willows, which, intermingling with peach and apricot
        trees, screened the heavens from view, and kept off the rays of the sun from
        this spot, which was in real truth devoid of even a grain of dust. Suddenly,
        they espied in the shade of the willows, an arched wooden bridge also reveal
        itself to the eye, with bannisters of vermilion colour.}
}
{
}

\Verse{73}
{
    \zA{}
}
{
    \eA{They crossed the bridge, and lo, all the paths lay open before them; but
        their gaze was readily attracted by a brick cottage spotless and
        cool-looking; whose walls were constructed of polished bricks, of uniform
        colour; (whose roof was laid) with speckless tiles; and whose enclosing
        walls were painted; while the minor slopes, which branched off from the main
        hill, all passed along under the walls on to the other side. "This house, in
        a site like this, is perfectly destitute of any charm!" added Chia Cheng.}
}
{
}

\Verse{74}
{
    \zA{}
}
{
    \eA{And as they entered the door, abruptly appeared facing them, a large boulder
        studded with holes and soaring high in the skies, which was surrounded on
        all four sides by rocks of every description, and completely, in fact, hid
        from view the rooms situated in the compound. But of flowers or trees, there
        was not even one about; and all that was visible were a few strange kinds of
        vegetation; some being of the creeper genus, others parasitic plants, either
        hanging from the apex of the hill, or inserting themselves into the base of
        the rocks; drooping down even from the eaves of the house, entwining the
        pillars, and closing round the stone steps. Or like green bands, they waved
        and flapped; or like gold thread, they coiled and bent, either with seeds
        resembling cinnabar, or with blossoms like golden olea;}
}
{
}

\Verse{75}
{
    \zA{}
}
{
    \eA{whose fragrance and aroma could not be equalled by those emitted by flowers
        of ordinary species. "This is pleasant!" Chia Cheng could not refrain from
        saying; "the only thing is that I don't know very much about flowers." "What
        are here are lianas and ficus pumila!" some of the company observed. "How
        ever can the liana and the ficus have such unusual scent?" questioned Chia
        Cheng. "Indeed they aren't!" interposed Pao-yü. "Among all these flowers,
        there are also ficus and liana, but those scented ones are iris, ligularia,
        and 'Wu' flowers; that kind consist, for the most part, of 'Ch'ih' flowers
        and orchids; while this mostly of gold-coloured dolichos. That species is
        the hypericum plant, this the 'Yü Lu' creeper. The red ones are, of course,
        the purple rue;}
}
{
}

\Verse{76}
{
    \zA{}
}
{
    \eA{the green ones consist for certain, of the green 'Chih' plant; and, to the
        best of my belief, these various plants are mentioned in the 'Li Sao' and
        'Wen Hsuan.' These rare plants are, some of them called something or other
        like 'Huo Na' and 'Chiang Hui;' others again are designated something like
        'Lun Tsu' and 'Tz'u Feng;' while others there are whose names sound like
        'Shih Fan,' 'Shui Sung' and 'Fu Liu,' which together with other species are
        to be found in the 'Treatise about the Wu city' by Tso T'ai-chung. There are
        also those which go under the appellation of 'Lu T'i,' or something like
        that; while there are others that are called something or other like 'Tan
        Chiao,' 'Mi Wu' and 'Feng Lien;'}
}
{
}

\Verse{77}
{
    \zA{}
}
{
    \eA{reference to which is made in the 'Treatise on the Shu city.' But so many
        years have now elapsed, and the times have so changed (since these treatises
        were written), that people, being unable to discriminate (the real names)
        may consequently have had to appropriate in every case such names as suited
        the external aspect, so that they may, it is quite possible, have gradually
        come to be called by wrong designations." But he had no time to conclude;
        for Chia Cheng interrupted him. "Who has ever asked you about it?" he
        shouted; which plunged Pao-yü into such a fright, that he drew back, and did
        not venture to utter another word.}
}
{
}

\Verse{78}
{
    \zA{}
}
{
    \eA{Chia Cheng perceiving that on both sides alike were covered passages
        resembling outstretched arms, forthwith continued his steps and entered the
        covered way, when he caught sight, at the upper end, of a five-roomed
        building, without spot or blemish, with folding blinds extending in a
        connected line, and with corridors on all four sides; (a building) which
        with its windows so green, and its painted walls, excelled, in spotless
        elegance, the other buildings they had seen before, to which it presented
        such a contrast. Chia Cheng heaved a sigh. "If one were able," he observed,
        "to boil his tea and thrum his lyre in here, there wouldn't even be any need
        for him to burn any more incense.}
}
{
}

\Verse{79}
{
    \zA{}
}
{
    \eA{But the execution of this structure is so beyond conception that you must,
        gentlemen, compose something nice and original to embellish the tablet with,
        so as not to render such a place of no effect!" "There's nothing so really
        pat," suggested the company smiling; "as 'the orchid-smell-laden breeze' and
        'the dew-bedecked epidendrum!" "These are indeed the only four characters,"
        rejoined Chia Cheng, "that could be suitably used; but what's to be said as
        far as the scroll goes?" "I've thought of a couplet," interposed one of the
        party, "which you'll all have to criticise, and put into ship-shape; its
        burden is this:  "The musk-like epidendrum smell enshrouds the court, where
        shines the  sun with oblique beams;}
}
{
}

\Verse{80}
{
    \zA{}
}
{
    \eA{The iris fragrance is wafted over the isle illumined by the moon's  clear
        rays." "As far as excellence is concerned, it's excellent," observed the
        whole party, "but the two words representing 'with oblique beams' are not
        felicitous." And as some one quoted the line from an old poem:  The angelica
        fills the court with tears, what time the sun doth slant. "Lugubrious,
        lugubrious!" expostulated the company with one voice. Another person then
        interposed. "I also have a couplet, whose merits you, gentlemen, can weigh;
        it runs as follows:  "Along the three pathways doth float the Yü Hui scented
        breeze! The radiant moon in the whole hall shines on the gold orchid!" Chia
        Cheng tugged at his moustache and gave way to meditation.}
}
{
}

\Verse{81}
{
    \zA{}
}
{
    \eA{He was just about also to suggest a stanza, when, upon suddenly raising his
        head, he espied Pao-yü standing by his side, too timid to give vent to a
        single sound. "How is it," he purposely exclaimed, "that when you should
        speak, you contrariwise don't? Is it likely that you expect some one to
        request you to confer upon us the favour of your instruction?" "In this
        place," Pao-yü rejoined at these words, "there are no such things as
        orchids, musk, resplendent moon or islands; and were one to begin quoting
        such specimens of allusions, to scenery, two hundred couplets could be
        readily given without, even then, having been able to exhaust the supply!"
        "Who presses your head down," Chia Cheng urged, "and uses force that you
        must come out with all these remarks?"}
}
{
}

\Verse{82}
{
    \zA{}
}
{
    \eA{"Well, in that case," added Pao-yü, "there are no fitter words to put on the
        tablet than the four representing: 'The fragrance pure of the ligularia and
        iris.' While the device on the scroll might be:  "Sung is the nutmeg song,
        but beauteous still is the sonnet! Near the T'u Mei to sleep, makes e'en a
        dream with fragrance full!" "This is," laughed Chia Cheng sneeringly, "an
        imitation of the line:  "A book when it is made of plaintain leaves, the
        writing green is also  bound to be! "So that there's nothing remarkable
        about it." "Li T'ai-po, in his work on the Phoenix Terrace," protested the
        whole party, "copied, in every point, the Huang Hua Lou. But what's
        essential is a faultless imitation.}
}
{
}

\Verse{83}
{
    \zA{}
}
{
    \eA{Now were we to begin to criticise minutely the couplet just cited, we would
        indeed find it to be, as compared with the line 'A book when it is made of
        plantain leaves,' still more elegant and of wider application!" "What an
        idea?" observed Chia Cheng derisively. But as he spoke, the whole party
        walked out; but they had not gone very far before they caught sight of a
        majestic summer house, towering high peak-like, and of a structure rising
        loftily with storey upon storey; and completely locked in as they were on
        every side they were as beautiful as the Jade palace. Far and wide, road
        upon road coiled and wound;}
}
{
}

\Verse{84}
{
    \zA{}
}
{
    \eA{while the green pines swept the eaves, the jady epidendrum encompassed the
        steps, the animals' faces glistened like gold, and the dragons' heads shone
        resplendent in their variegated hues. "This is the Main Hall," remarked Chia
        Cheng; "the only word against it is that there's a little too much finery."
        "It should be so," rejoined one and all, "so as to be what it's intended to
        be! The imperial consort has, it is true, an exalted preference for economy
        and frugality, but her present honourable position requires the observance
        of such courtesies, so that (finery) is no fault." As they made these
        remarks and advanced on their way the while, they perceived, just in front
        of them, an archway project to view, constructed of jadelike stone;}
}
{
}

\Verse{85}
{
    \zA{}
}
{
    \eA{at the top of which the coils of large dragons and the scales of small
        dragons were executed in perforated style. "What's the device to be for this
        spot?" inquired Chia Cheng. "It should be 'fairy land,'" suggested all of
        them, "so as to be apposite!" Chia Cheng nodded his head and said nothing.
        But as soon as Pao-yü caught sight of this spot something was suddenly
        aroused in his heart and he began to ponder within himself. "This place
        really resembles something that I've seen somewhere or other." But he could
        not at the moment recall to mind what year, moon, or day this had happened.
        Chia Cheng bade him again propose a motto;}
}
{
}

\Verse{86}
{
    \zA{}
}
{
    \eA{but Pao-yü was bent upon thinking over the details of the scenery he had
        seen on a former occasion, and gave no thought whatever to this place, so
        that the whole company were at a loss what construction to give to his
        silence, and came simply to the conclusion that, after the bullying he had
        had to put up with for ever so long, his spirits had completely vanished,
        his talents become exhausted and his speech impoverished; and that if he
        were harassed and pressed, he might perchance, as the result of anxiety,
        contract some ailment or other, which would of course not be a suitable
        issue, and they lost no time in combining together to dissuade Chia Cheng.
        "Never mind," they said, "to-morrow will do to compose some device; let's
        drop it now."}
}
{
}

\Verse{87}
{
    \zA{}
}
{
    \eA{Chia Cheng himself was inwardly afraid lest dowager lady Chia should be
        anxious, so that he hastily remarked as he forced a smile. "You beast, there
        are, after all, also occasions on which you are no good! but never mind!
        I'll give you one day to do it in, and if by to-morrow you haven't been able
        to compose anything, I shall certainly not let you off. This is the first
        and foremost place and you must exercise due care in what you write." Saying
        this, he sallied out, at the head of the company, and cast another glance at
        the scenery.}
}
{
}

\Verse{88}
{
    \zA{}
}
{
    \eA{Indeed from the time they had entered the gate up to this stage, they had
        just gone over five or six tenths of the whole ground, when it happened
        again that a servant came and reported that some one had arrived from Mr.
        Yü-'ts'un's to deliver a message. "These several places (which remain),"
        Chia Cheng observed with a smile, "we have no time to pass under inspection;
        but we might as well nevertheless go out at least by that way, as we shall
        be able, to a certain degree, to have a look at the general aspect." With
        these words, he showed the way for the family companions until they reached
        a large bridge, with water entering under it, looking like a curtain made of
        crystal.}
}
{
}

\Verse{89}
{
    \zA{}
}
{
    \eA{This bridge, the fact is, was the dam, which communicated with the river
        outside, and from which the stream was introduced into the grounds. "What's
        the name of this water-gate?" Chia Cheng inquired. "This is," replied
        Pao-yü, "the main stream of the Hsin Fang river, and is therefore called the
        Hsin Fang water-gate." "Nonsense!" exclaimed Chia Cheng. "The two words Hsin
        Fang must on no account be used!" And as they speedily advanced on their
        way, they either came across elegant halls, or thatched cottages; walls made
        of piled-up stone, or gates fashioned of twisted plants; either a secluded
        nunnery or Buddhist fane, at the foot of some hill; or some unsullied
        houses, hidden in a grove, tenanted by rationalistic priestesses; either
        extensive corridors and winding grottoes; or square buildings, and circular
        pavilions.}
}
{
}

\Verse{90}
{
    \zA{}
}
{
    \eA{But Chia Cheng had not the energy to enter any of these places, for as he
        had not had any rest for ever so long, his legs felt shaky and his feet
        weak. Suddenly they also discerned ahead of them a court disclose itself to
        view. "When we get there," Chia Cheng suggested, "we must have a little
        rest." Straightway as he uttered the remark, he led them in, and winding
        round the jade-green peach-trees, covered with blossom, they passed through
        the bamboo fence and flower-laden hedge, which were twisted in such a way as
        to form a circular, cavelike gateway, when unexpectedly appeared before
        their eyes an enclosure with whitewashed walls, in which verdant willows
        drooped in every direction.}
}
{
}

\Verse{91}
{
    \zA{}
}
{
    \eA{Chia Cheng entered the gateway in company with the whole party. Along the
        whole length of both sides extended covered passages, connected with each
        other; while in the court were laid out several rockeries. In one quarter
        were planted a number of banana trees; on the opposite stood a plant of
        begonia from Hsi Fu. Its appearance was like an open umbrella. The gossamer
        hanging (from its branches) resembled golden threads. The corollas (seemed)
        to spurt out cinnabar. "What a beautiful flower! what a beautiful flower!"
        ejaculated the whole party with one voice; "begonias are verily to be found;
        but never before have we seen anything the like of this in beauty." "This is
        called the maiden begonia and is, in fact, a foreign species," Chia Cheng
        observed.}
}
{
}

\Verse{92}
{
    \zA{}
}
{
    \eA{"There's a homely tradition that it is because it emanates from the maiden
        kingdom that its flowers are most prolific; but this is likewise erratic
        talk and devoid of common sense." "They are, after all," rejoined the whole
        company, "so unlike others (we have seen), that what's said about the maiden
        kingdom is, we are inclined to believe, possibly a fact." "I presume,"
        interposed Pao-yü, "that some clever bard or poet, (perceiving) that this
        flower was red like cosmetic, delicate as if propped up in sickness, and
        that it closely resembled the nature of a young lady, gave it, consequently,
        the name of maiden! People in the world will propagate idle tales, all of
        which are unavoidably treated as gospel!" "We receive (with thanks) your
        instructions;}
}
{
}

\Verse{93}
{
    \zA{}
}
{
    \eA{what excellent explanation!" they all remarked unanimously, and as they
        expressed these words, the whole company took their seats on the sofas under
        the colonnade. "Let's think of some original text or other for a motto,"
        Chia Cheng having suggested, one of the companions opined that the two
        characters: "Banana and stork" would be felicitous; while another one was of
        the idea that what would be faultless would be: "Collected splendour and
        waving elegance!" "'Collected splendour and waving elegance' is excellent,"
        Chia Cheng observed addressing himself to the party; and Pao-yü himself,
        while also extolling it as beautiful, went on to say: "There's only one
        thing however to be regretted!" "What about regret?" the company inquired.}
}
{
}

\Verse{94}
{
    \zA{}
}
{
    \eA{"In this place," Pao-yü explained, "are set out both bananas as well as
        begonias, with the intent of secretly combining in them the two properties
        of red and green; and if mention of one of them be made, and the other be
        omitted, (the device) won't be good enough for selection." "What would you
        then suggest?" Chia Cheng asked. "I would submit the four words, 'the red
        (flowers) are fragrant, the green (banana leaves) like jade,' which would
        render complete the beauties of both (the begonias and bananas)." "It isn't
        good! it isn't good!" Chia Cheng remonstrated as he shook his head;}
}
{
}

\Verse{95}
{
    \zA{}
}
{
    \eA{and while passing this remark, he conducted the party into the house, where
        they noticed that the internal arrangements effected differed from those in
        other places, as no partitions could, in fact, be discerned. Indeed, the
        four sides were all alike covered with boards carved hollow with fretwork,
        (in designs consisting) either of rolling clouds and hundreds of bats; or of
        the three friends of the cold season of the year, (fir, bamboo and almond);
        of scenery and human beings, or of birds or flowers; either of clusters of
        decoration, or of relics of olden times; either of ten thousand characters
        of happiness or of ten thousand characters of longevity.}
}
{
}

\Verse{96}
{
    \zA{}
}
{
    \eA{The various kinds of designs had been all carved by renowned hands, in
        variegated colours, inlaid with gold, and studded with precious gems; while
        on shelf upon shelf were either arranged collections of books, or tripods
        were laid out; either pens and inkslabs were distributed about, or vases
        with flowers set out, or figured pots were placed about; the designs of the
        shelves being either round or square; or similar to sunflowers or banana
        leaves; or like links, half overlapping each other. And in very truth they
        resembled bouquets of flowers or clusters of tapestry, with all their
        fretwork so transparent. Suddenly (the eye was struck) by variegated gauzes
        pasted (on the wood-work), actually forming small windows;}
}
{
}

\Verse{97}
{
    \zA{}
}
{
    \eA{and of a sudden by fine thin silks lightly overshadowing (the fretwork) just
        as if there were, after all, secret doors. The whole walls were in addition
        traced, with no regard to symmetry, with outlines of the shapes of curios
        and nick-nacks in imitation of lutes, double-edged swords, hanging bottles
        and the like, the whole number of which, though (apparently) suspended on
        the walls, were all however on a same level with the surface of the
        partition walls. "What fine ingenuity!" they all exclaimed extollingly;}
}
{
}

\Verse{98}
{
    \zA{}
}
{
    \eA{"what a labour they must have been to carry out!" Chia Cheng had actually
        stepped in; but scarcely had they reached the second stage, before the whole
        party readily lost sight of the way by which they had come in. They glanced
        on the left, and there stood a door, through which they could go. They cast
        their eyes on the right, and there was a window which suddenly impeded their
        progress. They went forward, but there again they were obstructed by a
        bookcase. They turned their heads round, and there too stood windows pasted
        with transparent gauze and available door-ways: but the moment they came
        face to face with the door, they unexpectedly perceived that a whole company
        of people had likewise walked in, just in front of them, whose appearance
        resembled their own in every respect.}
}
{
}

\Verse{99}
{
    \zA{}
}
{
    \eA{But it was only a mirror. And when they rounded the mirror, they detected a
        still larger number of doors. "Sir," Chia Chen remarked with a grin; "if
        you'll follow me out through this door, we'll forthwith get into the
        back-court; and once out of the back-court, we shall be, at all events,
        nearer than we were before." Taking the lead, he conducted Chia Cheng and
        the whole party round two gauze mosquito houses, when they verily espied a
        door through which they made their exit, into a court, replete with stands
        of cinnamon roses. Passing round the flower-laden hedge, the only thing that
        spread before their view was a pure stream impeding their advance. The whole
        company was lost in admiration.}
}
{
}

\Verse{100}
{
    \zA{}
}
{
    \eA{"Where does this water again issue from?" they cried. Chia Chen pointed to a
        spot at a distance. "Starting originally," he explained, "from that
        water-gate, it runs as far as the mouth of that cave, when from among the
        hills on the north-east side, it is introduced into that village, where
        again a diverging channel has been opened and it is made to flow in a
        south-westerly direction; the whole volume of water then runs to this spot,
        where collecting once more in one place, it issues, on its outward course,
        from beneath that wall." "It's most ingenious!" they one and all exclaimed,
        after they had listened to him; but, as they uttered these words, they
        unawares realised that a lofty hill obstructed any further progress. The
        whole party felt very hazy about the right road.}
}
{
}

\Verse{101}
{
    \zA{}
}
{
    \eA{But "Come along after me," Chia Chen smilingly urged, as he at once went
        ahead and showed the way, whereupon the company followed in his steps, and
        as soon as they turned round the foot of the hill, a level place and broad
        road lay before them; and wide before their faces appeared the main
        entrance. "This is charming! this is delightful!" the party unanimously
        exclaimed, "what wits must have been ransacked, and ingenuity attained, so
        as to bring things to this extreme degree of excellence!" Forthwith the
        party egressed from the garden, and Pao-yü's heart anxiously longed for the
        society of the young ladies in the inner quarters, but as he did not hear
        Chia Cheng bid him go, he had no help but to follow him into the library.}
}
{
}

\Verse{102}
{
    \zA{}
}
{
    \eA{But suddenly Chia Cheng bethought himself of him. "What," he said, "you
        haven't gone yet! the old lady will I fear be anxious on your account; and
        is it pray that you haven't as yet had enough walking?" Pao-yü at length
        withdrew out of the library. On his arrival in the court, a page, who had
        been in attendance on Chia Cheng, at once pressed forward, and took hold of
        him fast in his arms. "You've been lucky enough," he said, "to-day to have
        been in master's good graces! just a while back when our old mistress
        despatched servants to come on several occasions and ask after you, we
        replied that master was pleased with you; for had we given any other answer,
        her ladyship would have sent to fetch you to go in, and you wouldn't have
        had an opportunity of displaying your talents.}
}
{
}

\Verse{103}
{
    \zA{}
}
{
    \eA{Every one admits that the several stanzas you recently composed were
        superior to those of the whole company put together; but you must, after the
        good luck you've had to-day, give us a tip!" "I'll give each one of you a
        tiao," Pao-yü rejoined smirkingly. "Who of us hasn't seen a tiao?" they all
        exclaimed, "let's have that purse of yours, and have done with it!" Saying
        this, one by one advanced and proceeded to unloosen the purse, and to
        unclasp the fan-case; and allowing Pao-yü no time to make any remonstrance,
        they stripped him of every ornament in the way of appendage which he carried
        about on his person. "Whatever we do let's escort him home!" they shouted,
        and one after another hustled round him and accompanied him as far as
        dowager lady Chia's door.}
}
{
}

\Verse{104}
{
    \zA{}
}
{
    \eA{Her ladyship was at this moment awaiting his arrival, so that when she saw
        him walk in, and she found out that (Chia Cheng) had not bullied him, she
        felt, of course, extremely delighted. But not a long interval elapsed before
        Hsi Jen came to serve the tea; and when she perceived that on his person not
        one of the ornaments remained, she consequently smiled and inquired: "Have
        all the things that you had on you been again taken away by these barefaced
        rascals?" As soon as Lin Tai-yü heard this remark, she crossed over to him
        and saw at a glance that not one single trinket was, in fact, left.}
}
{
}

\Verse{105}
{
    \zA{}
}
{
    \eA{"Have you also given them," she felt constrained to ask, "the purse that I
        gave you? Well, by and by, when you again covet anything of mine, I shan't
        let you have it." After uttering these words, she returned into her
        apartment in high dudgeon, and taking the scented bag, which Pao-yü had
        asked her to make for him, and which she had not as yet finished, she picked
        up a pair of scissors, and instantly cut it to pieces. Pao-yü noticing that
        she had lost her temper, came after her with hurried step, but the bag had
        already been cut with the scissors;}
}
{
}

\Verse{106}
{
    \zA{}
}
{
    \eA{and as Pao-yü observed how extremely fine and artistic this scented bag was,
        in spite of its unfinished state, he verily deplored that it should have
        been rent to pieces for no rhyme or reason. Promptly therefore unbuttoning
        his coat, he produced from inside the lapel the purse, which had been
        fastened there. "Look at this!" he remarked as he handed it to Tai-yü; "what
        kind of thing is this! have I given away to any one what was yours?" Lin
        Tai-yü, upon seeing how much he prized it as to wear it within his clothes,
        became alive to the fact that it was done with intent, as he feared lest any
        one should take it away; and as this conviction made her sorry that she had
        been so impetuous as to have cut the scented bag, she lowered her head and
        uttered not a word.}
}
{
}

\Verse{107}
{
    \zA{}
}
{
    \eA{"There was really no need for you to have cut it," Pao-yü observed; "but as
        I know that you're loth to give me anything, what do you say to my returning
        even this purse?" With these words, he threw the purse in her lap and walked
        off; which vexed Tai-yü so much the more that, after giving way to tears,
        she took up the purse in her hands to also destroy it with the scissors,
        when Pao-yü precipitately turned round and snatched it from her grasp. "My
        dear cousin," he smilingly pleaded, "do spare it!" and as Tai-yü dashed down
        the scissors and wiped her tears: "You needn't," she urged, "be kind to me
        at one moment, and unkind at another;}
}
{
}

\Verse{108}
{
    \zA{}
}
{
    \eA{if you wish to have a tiff, why then let's part company!" But as she spoke,
        she lost control over her temper, and, jumping on her bed, she lay with her
        face turned towards the inside, and set to work drying her eyes. Pao-yü
        could not refrain from approaching her. "My dear cousin, my own cousin," he
        added, "I confess my fault!" "Go and find Pao-yü!" dowager lady Chia
        thereupon gave a shout from where she was in the front apartment, and all
        the attendants explained that he was in Miss Lin's room. "All right, that
        will do! that will do!" her ladyship rejoined, when she heard this reply;
        "let the two cousins play together; his father kept him a short while back
        under check, for ever so long, so let him have some distraction.}
}
{
}

\Verse{109}
{
    \zA{}
}
{
    \eA{But the only thing is that you mustn't allow them to have any quarrels." To
        which the servants in a body expressed their obedience. Tai-yü, unable to
        put up with Pao-yü's importunity, felt compelled to rise. "Your object seems
        to be," she remarked, "not to let me have any rest. If it is, I'll run away
        from you." Saying which, she there and then was making her way out, when
        Pao-yü protested with a face full of smiles: "Wherever you go, I'll follow!"
        and as he, at the same time, took the purse and began to fasten it on him,
        Tai-yü stretched out her hand, and snatching it away, "You say you don't
        want it," she observed, "and now you put it on again! I'm really much
        ashamed on your account!"}
}
{
}

\Verse{110}
{
    \zA{}
}
{
    \eA{And these words were still on her lips when with a sound of Ch'ih, she burst
        out laughing. "My dear cousin," Pao-yü added, "to-morrow do work another
        scented bag for me!" "That too will rest upon my good pleasure," Tai-yü
        rejoined. As they conversed, they both left the room together and walked
        into madame Wang's suite of apartments, where, as luck would have it,
        Pao-ch'ai was also seated. Unusual commotion prevailed, at this time, over
        at madame Wang's, for the fact is that Chia Se had already come back from Ku
        Su, where he had selected twelve young girls, and settled about an
        instructor, as well as about the theatrical properties and the other
        necessaries.}
}
{
}

\Verse{111}
{
    \zA{}
}
{
    \eA{And as Mrs. Hsüeh had by this date moved her quarters into a separate place
        on the northeast side, and taken up her abode in a secluded and quiet house,
        (madame Wang) had had repairs of a distinct character executed in the Pear
        Fragrance Court, and then issued directions that the instructor should train
        the young actresses in this place; and casting her choice upon all the
        women, who had, in days of old, received a training in singing, and who were
        now old matrons with white hair, she bade them have an eye over them and
        keep them in order. Which done, she enjoined Chia Se to assume the chief
        control of all matters connected with the daily and monthly income and
        outlay, as well as of the accounts of all articles in use of every kind and
        size.}
}
{
}

\Verse{112}
{
    \zA{}
}
{
    \eA{Lin Chih-hsiao also came to report: "that the twelve young nuns and Taoist
        girls, who had been purchased after proper selection, had all arrived, and
        that the twenty newly-made Taoist coats had also been received. That there
        was besides a maiden, who though devoted to asceticism, kept her chevelure
        unshaved; that she was originally a denizen of Suchow, of a family whose
        ancestors were also people of letters and official status; that as from her
        youth up she had been stricken with much sickness, (her parents) had
        purchased a good number of substitutes (to enter the convent), but all with
        no relief to her, until at last this girl herself entered the gate of
        abstraction when she at once recovered.}
}
{
}

\Verse{113}
{
    \zA{}
}
{
    \eA{That hence it was that she grew her hair, while she devoted herself to an
        ascetic life; that she was this year eighteen years of age, and that the
        name given to her was Miao Yü; that her father and mother were, at this
        time, already dead; that she had only by her side, two old nurses and a
        young servant girl to wait upon her; that she was most proficient in
        literature, and exceedingly well versed in the classics and canons; and that
        she was likewise very attractive as far as looks went; that having heard
        that in the city of Ch'ang-an, there were vestiges of Kuan Yin and relics of
        the canons inscribed on leaves, she followed, last year, her teacher (to the
        capital). She now lives," he said, "in the Lao Ni nunnery, outside the
        western gate;}
}
{
}

\Verse{114}
{
    \zA{}
}
{
    \eA{her teacher was a great expert in prophetic divination, but she died in the
        winter of last year, and her dying words were that as it was not suitable
        for (Miao Yü) to return to her native place, she should await here, as
        something in the way of a denouement was certain to turn up; and this is the
        reason why she hasn't as yet borne the coffin back to her home!" "If such be
        the case," madame Wang readily suggested, "why shouldn't we bring her here?"
        "If we are to ask her," Lin Chih-hsiao's wife replied, "she'll say that a
        marquis' family and a duke's household are sure, in their honourable
        position, to be overbearing to people; and I had rather not go." "As she's
        the daughter of an official family," madame Wang continued, "she's bound to
        be inclined to be somewhat proud;}
}
{
}

\Verse{115}
{
    \zA{}
}
{
    \eA{but what harm is there to our sending her a written invitation to ask her to
        come!" Lin Chih-hsiao's wife assented; and leaving the room, she made the
        secretary write an invitation and then went to ask Miao Yü. The next day
        servants were despatched, and carriages and sedan chairs were got ready to
        go and bring her over. What subsequently transpired is not as yet known,
        but, reader, listen to the account given in the following chapter.}
}
{
}
